\documentclass[aspectratio=169,10pt,t]{beamer}
\usepackage{ac1}
\aulainfo{10}{Subalgoritmos}{Funções, fluxo de execução e return}

\begin{document}
\capa

\begin{frame}{Roteiro}
  \begin{itemize}
    \item Por que dividir um programa em partes? (Cap.~4)
    \item Fluxo de execução (4.1)
    \item Definindo e chamando funções (4.3)
    \item Funções sem retorno: procedimentos
    \item Funções com retorno: \texttt{return}
    \item Documentação com \emph{docstrings}
  \end{itemize}
\end{frame}

\begin{frame}{Dividir para conquistar}
  \begin{columns}[T]
    \column{0.55\textwidth}
    \begin{itemize}
      \item Um problema grande vira \textbf{vários problemas pequenos}
      \item Cada parte é um \textbf{subalgoritmo}: em Python, uma \alert{função}
      \item Benefícios:
      \begin{itemize}\small
        \item \textbf{reuso}: escreve uma vez, usa várias
        \item \textbf{legibilidade}: nomes explicam o que cada parte faz
        \item \textbf{teste}: verifica-se cada função isoladamente
      \end{itemize}
    \end{itemize}
    \column{0.42\textwidth}
    \begin{exampleblock}{Já usamos funções!}
      \texttt{print()}, \texttt{input()}, \texttt{int()}, \texttt{range()}, \texttt{len()}, \texttt{type()}\ldots\\
      Agora vamos \emph{criar} as nossas.
    \end{exampleblock}
  \end{columns}
\end{frame}

\begin{frame}[fragile]{Definindo e chamando uma função}
\begin{lstlisting}
def imprime_linha():
    print('**********')

imprime_linha()
imprime_linha()
\end{lstlisting}
\begin{lstlisting}[style=saida]
**********
**********
\end{lstlisting}
  \begin{itemize}
    \item \texttt{def nome(parâmetros):} \(\to\) \alert{define} (apenas guarda; ainda não executa)
    \item \texttt{nome()} \(\to\) \alert{chama} (executa o corpo)
    \item O corpo é um bloco indentado
  \end{itemize}
  \fonte{code/Cap4/prog.4.8\_imprimeb.py}
\end{frame}

\begin{frame}[fragile]{Fluxo de execução}
  \begin{columns}[T]
    \column{0.5\textwidth}
\begin{lstlisting}
def a():
    print("início de a")
    b()
    print("fim de a")

def b():
    print("em b")

print("começo")
a()
print("fim")
\end{lstlisting}
    \column{0.45\textwidth}
    \textbf{Ordem de execução}
    \begin{enumerate}\small
      \item as definições \texttt{def} são lidas (nada executa)
      \item \texttt{print("começo")}
      \item chama \texttt{a()} \(\to\) pula para o corpo
      \item chama \texttt{b()} \(\to\) pula de novo
      \item \texttt{b} termina \(\to\) volta a \texttt{a}
      \item \texttt{a} termina \(\to\) volta ao programa
      \item \texttt{print("fim")}
    \end{enumerate}
  \end{columns}
\end{frame}

\begin{frame}[fragile]{Funções com parâmetros}
\begin{lstlisting}
def imprime_linha(n):
    print('*' * n)

imprime_linha(5)
imprime_linha(20)
\end{lstlisting}
\begin{lstlisting}[style=saida]
*****
********************
\end{lstlisting}
  \begin{itemize}
    \item \textbf{Parâmetro}: nome na definição (\texttt{n}); \textbf{argumento}: valor na chamada (\texttt{5})
    \item \texttt{'*' * n} repete a string \(n\) vezes
    \item Exercício 4.5: usar funções para desenhar outras figuras com asteriscos
  \end{itemize}
  \fonte{code/Cap4/prog.4.9\_imprimenb.py}
\end{frame}

\begin{frame}[fragile]{Funções sem retorno (procedimentos)}
\begin{lstlisting}
def conceito(nota):
    """Imprime o conceito correspondente à nota."""
    if nota >= 9:
        print("A")
    elif nota >= 7:
        print("B")
    elif nota >= 5:
        print("C")
    else:
        print("D")

conceito(8.5)
\end{lstlisting}
  \begin{itemize}
    \item A função \alert{faz algo} (imprime) mas não devolve valor ao chamador
    \item Exercício 4.7 do livro (nota \(\to\) conceito)
  \end{itemize}
\end{frame}

\begin{frame}[fragile]{Funções com retorno: \texttt{return}}
\begin{lstlisting}
def quadrado(x):
    return x * x

a = quadrado(7)
print(a)                   # 49
print(quadrado(3) + 1)     # 10
\end{lstlisting}
  \begin{alertblock}{\texttt{print} $\neq$ \texttt{return}}
    \texttt{print} \emph{mostra} na tela; \texttt{return} \emph{devolve} o valor para quem chamou, que pode usá-lo em contas. Ao executar \texttt{return}, a função termina.
  \end{alertblock}
  Exercício 4.8: função que recebe um inteiro e devolve seu quadrado.
\end{frame}

\begin{frame}[fragile]{Funções matemáticas: conversão de temperaturas}
\begin{lstlisting}
def fahrenheit(celsius):
    """Converte Celsius para Fahrenheit."""
    return 9 / 5 * celsius + 32

def celsius(fahrenheit):
    """Converte Fahrenheit para Celsius."""
    return 5 / 9 * (fahrenheit - 32)

print(fahrenheit(100))     # 212.0
print(celsius(212))        # 100.0
\end{lstlisting}
  \begin{itemize}
    \item Exercícios 4.9 e 4.10; a ``ida e volta'' é um bom teste: \texttt{celsius(fahrenheit(x))} deve devolver \texttt{x}
  \end{itemize}
\end{frame}

\begin{frame}[fragile]{Fatorial como função}
\begin{lstlisting}
def fatorial(n):
    """Calcula o fatorial de n."""
    fat = 1
    for i in range(1, n + 1):
        fat = fat * i
    return fat

n = int(input('Qual o número para calcular o fatorial? '))
print('O fatorial de', n, 'é', fatorial(n))
\end{lstlisting}
  \begin{itemize}
    \item Compare com o programa sem função (Aula 8): a lógica foi \emph{empacotada} com um nome
    \item Com a função pronta, calcular \(\binom{n}{k}=\dfrac{n!}{k!(n-k)!}\) é trivial
  \end{itemize}
  \fonte{code/Cap4/prog.4.3\_func\_fatorial2.py}
\end{frame}

\begin{frame}[fragile]{Função com várias decisões e retorno}
\begin{lstlisting}
import math

def distancia(x1, y1, x2, y2):
    """Distância entre os pontos (x1, y1) e (x2, y2)."""
    return math.sqrt((x2 - x1)**2 + (y2 - y1)**2)

print(distancia(0, 0, 3, 4))      # 5.0
\end{lstlisting}
  \begin{itemize}
    \item Exercício 4.11: implemente e teste com o triângulo 3--4--5
    \item Funções podem chamar outras funções (inclusive de módulos importados)
  \end{itemize}
\end{frame}

\begin{frame}[fragile]{Documentando: docstrings}
\begin{lstlisting}
def area_circulo(raio):
    """Devolve a área do círculo de raio dado.

    raio -- número não negativo, em qualquer unidade.
    """
    return 3.14159 * raio ** 2

help(area_circulo)
\end{lstlisting}
  \begin{itemize}
    \item A primeira string do corpo é a \textbf{docstring}: documenta \emph{o que} a função faz
    \item \texttt{help()} a mostra; boa prática desde o Capítulo 9 do livro
    \item Nomes de funções: verbos em minúsculas com \texttt{\_} (\texttt{calcula\_media})
  \end{itemize}
\end{frame}

\begin{frame}{Atividade: de um programa longo a funções}
  Reescreva o programa da Aula 9 (média com validação) em três funções:
  \vspace{1ex}
  \begin{enumerate}
    \item \texttt{le\_nota(i)}: pede a nota \(i\) até ser válida e a devolve
    \item \texttt{calcula\_media(soma, n)}: devolve a média
    \item \texttt{situacao(media)}: devolve \texttt{"Aprovado"}, \texttt{"Prova final"} ou \texttt{"Reprovado"}
  \end{enumerate}
  \vspace{1ex}
  \begin{block}{Critério de qualidade}
    O programa principal deve ter menos de 10 linhas e ser \emph{legível como uma receita}.
  \end{block}
\end{frame}

\exercicios{4}{%
  \begin{itemize}\small
    \item Funções sem retorno: \textbf{4.5} (figura com asteriscos), \textbf{4.6} (árvore de Natal com \texttt{for}), \textbf{4.7} (conceito)
    \item Funções com retorno: \textbf{4.8} (quadrado), \textbf{4.9} (Fahrenheit), \textbf{4.10} (Celsius), \textbf{4.11} (distância)
  \end{itemize}}

\resumo{
  \item Função = bloco nomeado e reutilizável: \texttt{def nome(parâmetros):}
  \item Definir não executa; \emph{chamar} executa e depois o fluxo retorna
  \item \texttt{return} devolve um valor; \texttt{print} apenas imprime
  \item Docstrings documentam a função
  \item Próxima aula: parâmetros, escopo e múltiplos retornos
}

\end{document}
